% ==============================================================================
% Dados do trabalho
% Preencha os campos abaixo. Eles alimentam a capa, a folha de rosto, a folha
% de aprovação e as informações do PDF.
% ==============================================================================

% Curso: cc (Ciência da Computação) ou ec (Engenharia de Computação).
% Este exemplo usa ec apenas para mostrar a opção.
\curso{ec}

% Título e subtítulo (deixe \subtitulo{} vazio se não houver).
\titulo{Monitora}
\subtitulo{um sistema web para agendamento de monitorias acadêmicas}

\autor{Nome Completo do Autor}
\data{2026}

% Orientador e coorientador, com titulação. Para orientadora, use
% \orientador[Orientadora:]{Profa. Dra. Nome}. Sem coorientador, apague a linha.
\orientador{Prof. Dr. Nome do Orientador}
% \coorientador{Prof. Me. Nome do Coorientador}

% Data da defesa. Se ficar vazio, a folha de aprovação sai com linhas em branco.
\dataaprovacao{}

% Banca examinadora: \membrobanca{Nome com titulação}{Instituição}{Função}.
% O orientador é o primeiro membro.
\membrobanca{Prof. Dr. Nome do Orientador}{Universidade Federal do Espírito Santo}{Orientador}
\membrobanca{Prof. Dr. Nome do Examinador}{Universidade Federal do Espírito Santo}{Examinador interno}
\membrobanca{Profa. Dra. Nome da Examinadora}{Instituição de origem}{Examinadora externa}

% Informações do arquivo PDF.
\makeatletter
\hypersetup{
  pdftitle   = {\@title},
  pdfauthor  = {\@author},
  pdfsubject = {\imprimirtipotrabalho},
  pdfkeywords= {sistema web; agendamento; monitoria; engenharia de software},
}
\makeatother
